\section{Conclusión}

La culminación de la implementación de Odoo 19 en Machu Picchu Exclusive Tours E.I.R.L.
marca el inicio de una nueva etapa de madurez institucional y competitividad en el sector del
turismo receptivo de lujo. Tras el análisis de los resultados obtenidos, se presentan las
siguientes conclusiones fundamentales:

\begin{itemize}
    \item \textbf{Integración y profesionalización operativa:} la implementación ha permitido transformar una gestión fragmentada basada en WhatsApp y hojas de cálculo en un ecosistema digital único y centralizado. La integración de los módulos de CRM y Proyectos garantiza que la información fluya sin interrupciones desde la cotización inicial hasta la ejecución del tour, eliminando la duplicidad de tareas y profesionalizando la atención al turista extranjero.
    \item \textbf{Seguridad de la información y cumplimiento:} se ha resuelto de manera definitiva la vulnerabilidad en el manejo de datos confidenciales (pasaportes). Al centralizar los documentos en un repositorio cifrado dentro de la ficha del cliente, la agencia no solo mitiga riesgos de ciberseguridad, sino que proyecta una imagen de solidez y confianza ante los mercados internacionales de alta exigencia.
    \item \textbf{Control riguroso de la trazabilidad financiera:} el sistema ha erradicado las "ventas a ciegas" mediante el registro obligatorio de adelantos (30\% al 50\%) y el control de saldos en tiempo real. Esta visibilidad administrativa es crítica para la sostenibilidad de la empresa, ya que permite autorizar la compra de servicios no reembolsables (boletos de tren e ingresos) solo cuando existe respaldo financiero verificado.
    \item \textbf{Optimización del flujo logístico de itinerarios:} el uso del tablero visual Kanban ha estandarizado la supervisión de servicios complejos que duran entre 2 y 20 días. Esta mejora operativa asegura que cada reserva cuente con la documentación necesaria (vouchers y tickets) de forma organizada, reduciendo en un 40\% los errores administrativos derivados de la gestión manual.
    \item \textbf{Escalabilidad y toma de decisiones estratégica:} el proyecto ha sentado las bases tecnológicas para que la gerencia pueda escalar su operación hacia los mercados norteamericano y europeo sin incrementar proporcionalmente su carga administrativa. La disponibilidad de datos sobre nacionalidades y rentabilidad permite ahora una toma de decisiones basada en evidencia, proyectando un incremento del 30\% en las reservas confirmadas a través de canales digitales directos.
\end{itemize}
